\chapter{Introduction}

\section{Background and Motivation}

Healthcare fraud costs the US Medicare and Medicaid programmes over \$300 billion annually. The rapid growth of electronic claims databases and advances in machine learning (ML) and Explainable AI (XAI) now enable automated, interpretable detection of fraud patterns at scale. In the Indian context, the NFHS-5 survey (724,115 female respondents, 2019--2021) provides critical population health data for national policy; logical inconsistencies in self-reported responses must be identified to preserve its epidemiological reliability. This project addresses both challenges within a unified, investigator-facing decision-support system.

\section{Problem Statement}

\begin{itemize}
    \item \textbf{Provider-Level Fraud:} Among 5,410 CMS SynPUF Medicare providers, 506 (9.35\%) are labelled fraudulent; severe class imbalance (1:9.7) and interpretability requirements challenge standard classifiers.
    \item \textbf{Claim-Level Screening:} Individual claims must be assessed for transparent, field-traceable indicators — anomalous reimbursements, duplicate billing, temporal inconsistencies, and high-frequency patterns.
    \item \textbf{NFHS-5 Clinical Consistency:} Among 724,115 NFHS-5 respondents, clinically implausible combinations (e.g., blood glucose $\geq$200~mg/dL without a diabetes diagnosis) must be identified and risk-stratified at scale.
    \item \textbf{Operational Deployment:} Investigators need integrated claim search, provider profiling, beneficiary history, and case management in a single application with sub-10~ms query latency.
\end{itemize}

\section{Existing Systems and Their Limitations}

\begin{itemize}
    \item \textbf{Rule-Based Expert Systems:} Cannot adapt to evolving fraud tactics or detect aggregate anomalies distributed across multiple claims.
    \item \textbf{Statistical Anomaly Detection:} High false-positive rates; cannot distinguish legitimate outliers from genuine fraud.
    \item \textbf{Supervised ML Without Explainability:} High AUC but opaque scores; investigators cannot justify flags for audit or legal proceedings.
    \item \textbf{Single-Granularity and No Clinical Integration:} Existing tools operate at one level only and none integrates Medicare fraud detection with clinical survey validation.
\end{itemize}

\section{Need for the Proposed System}

\begin{itemize}
    \item \textbf{High-Performance Query Engine:} 558,211 claims and 724,115 NFHS-5 records demand sub-10~ms analytical query latency; DuckDB's embedded columnar engine satisfies this without a server process.
    \item \textbf{Interpretable Risk Scores:} SHAP TreeExplainer values provide per-provider feature-level explanations, enabling investigators to document evidence and adjust thresholds dynamically.
    \item \textbf{Persistent Case Management:} An SQLite-backed Investigation Workspace provides case creation, status tracking, note annotation, and audit-trail retrieval within the same interface.
\end{itemize}

\section{Project Objectives}

\begin{enumerate}
    \item Design and implement a data engineering pipeline ingesting raw CMS SynPUF CSV files and constructing an 85~MB DuckDB analytical database with sub-10~ms query latency.
    \item Develop and evaluate Random Forest, Logistic Regression, and XGBoost classifiers under stratified five-fold cross-validation with six evaluation metrics.
    \item Implement a SHAP TreeExplainer module for per-provider waterfall plots and global feature-importance rankings.
    \item Design six weighted clinical inconsistency rules for NFHS-5 validation and apply them to all 724,115 records with five-tier risk categorisation.
    \item Implement five transparent claim-level screening rules covering reimbursement outliers, duplicate billing, duration anomalies, complex procedures, and high-frequency beneficiaries.
    \item Build a nine-page Streamlit dashboard integrating all analytical components into a single browser-accessible application.
    \item Implement a configurable dual-threshold risk classification scheme (default: 0.30 Medium, 0.70 High) with real-time slider controls.
    \item Validate all core modules via a pytest test suite covering data loading, feature engineering, model prediction, claim screening, and investigation workspace operations.
\end{enumerate}

\section{Project Scope}

\begin{itemize}
    \item \textbf{Datasets:} CMS SynPUF Medicare (558,211 claims, 5,410 providers) and NFHS-5 (724,115 records); supervised provider-level fraud detection and rule-based claim and clinical validation.
    \item \textbf{Platform:} Fully local Streamlit application backed by DuckDB and SQLite; no cloud deployment, live claims adjudication, or NLP on unstructured text.
    \item \textbf{Data:} Exclusively synthetic SynPUF and publicly released NFHS-5 data; no identifiable patient information is processed.
\end{itemize}

\section{Proposed Solution Overview}

The proposed system is a Streamlit-based nine-page interactive decision-support dashboard. The first component applies six weighted clinical inconsistency rules to 724,115 NFHS-5 records, assigning each a composite risk score mapping to five tiers (Clean, Low, Medium, High, Critical). The second component engineers ten provider-level aggregate features from CMS SynPUF data and trains three classifiers — Random Forest, Logistic Regression, and XGBoost — under five-fold stratified cross-validation; XGBoost achieves the highest ROC-AUC (93.55\%) and F1 (57.56\%), while Random Forest achieves the highest precision (68.93\%). SHAP TreeExplainer provides per-provider feature attribution on demand. All data are served via a DuckDB database (85~MB, sub-10~ms latency), and investigation cases are persisted in a SQLite-backed workspace with a six-status lifecycle.

\section{Key Contributions}

\begin{itemize}
    \item \textbf{Dual-Domain Analytical Framework:} First unified platform combining US Medicare ML fraud detection with Indian NFHS-5 clinical data consistency validation.
    \item \textbf{Ten-Feature Provider Engineering Pipeline:} Reproducible end-to-end pipeline deriving ten provider-level features from raw CMS SynPUF CSV files.
    \item \textbf{Weighted Clinical Rule System:} Six-rule framework applied to 724,115 NFHS-5 records, identifying 20,514 (2.83\%) clinically implausible records across five risk tiers.
    \item \textbf{SHAP-Augmented Provider Risk Profiling:} On-demand, query-driven SHAP waterfall explanations without offline batch reprocessing.
    \item \textbf{Transparent Claim-Level Screening:} Five deterministic, field-traceable indicators covering reimbursement outliers, duplicate billing, duration anomalies, procedure complexity, and beneficiary frequency.
    \item \textbf{DuckDB High-Performance Backend and SQLite Case Management:} Sub-10~ms columnar query engine over 558,211 claims with a full investigation lifecycle manager.
\end{itemize}

\section{Organisation of the Report}

\begin{center}
\begin{tabular}{|l|l|}
\hline
\textbf{Chapter} & \textbf{Content} \\
\hline
2 — Literature Survey      & Review of fraud detection, anomaly detection, XAI, and clinical data quality methods \\
3 — System Design          & Architecture, data-flow diagrams, ER schemas, and design rationale \\
4 — Implementation         & Module-level technical descriptions, code excerpts, and configuration details \\
5 — Results and Discussion & ML metrics, NFHS-5 flagging statistics, SHAP importances, and query benchmarks \\
6 — Conclusion             & Achievements, limitations, future enhancements, and deployment scenarios \\
7 — Research Publication   & Research paper prepared for peer-reviewed conference or journal submission \\
\hline
\end{tabular}
\end{center}
